% ===== PRÉAMBULE CANONIQUE — à coller VERBATIM en tête de CHAQUE .tex =====
\documentclass[11pt,a4paper]{article}
\usepackage[utf8]{inputenc}\usepackage[T1]{fontenc}\usepackage{lmodern}
\usepackage[french]{babel}
\usepackage{amsmath,amssymb,mathtools}\usepackage{icomma}
\usepackage{siunitx}\sisetup{locale=FR}  % \qty{}{}, \si{}, \ang{}
\usepackage{xcolor}\usepackage{colortbl}
\usepackage{tikz}\usepackage{pgfplots}\pgfplotsset{compat=1.18}
\usepackage[siunitx,europeanresistors]{circuitikz} % circuits (élec.)
\usepackage[most]{tcolorbox}\usepackage{enumitem}\usepackage{array,booktabs}
\usepackage[a4paper,margin=2.2cm]{geometry}
\usetikzlibrary{arrows.meta,calc,angles,quotes,positioning,patterns,%
  backgrounds,shapes.geometric,decorations.pathmorphing,decorations.markings,intersections}
\definecolor{primary}{RGB}{222,49,99}\definecolor{primarydark}{RGB}{150,22,55}
\definecolor{defcol}{RGB}{0,114,178}\definecolor{methcol}{RGB}{52,92,148}
\definecolor{excol}{RGB}{0,135,90}\definecolor{attncol}{RGB}{200,40,55}
\tcbset{boxbase/.style={enhanced,breakable,boxrule=0.9pt,arc=2.5pt,
  left=3mm,right=3mm,top=2.3mm,bottom=2.3mm,fonttitle=\bfseries,coltitle=white}}
% --- boîtes TUTORIEL (noms EXACTS, ne pas renommer) ---
\newtcolorbox{cle}[1][]{boxbase,colback=defcol!6,colframe=defcol,
  title={\textsc{À retenir}\ifx\relax#1\relax\relax\else\textmd{ : \itshape #1}\fi}}
\newtcolorbox{methode}[1][]{boxbase,colback=methcol!7,colframe=methcol,sharp corners,
  title={\textsc{Méthode}\ifx\relax#1\relax\relax\else\textmd{ --- \itshape #1}\fi}}
\newtcolorbox{exemplebox}[1][]{boxbase,colback=excol!5,colframe=excol,
  title={\textsc{Exemple}\ifx\relax#1\relax\relax\else\textmd{ : \itshape #1}\fi}}
\newtcolorbox{piege}[1][]{boxbase,colback=attncol!6,colframe=attncol,
  title={\textsc{Piège}\ifx\relax#1\relax\relax\else\textmd{ : \itshape #1}\fi}}
\newtcolorbox{remarque}[1][]{boxbase,colback=black!4,colframe=black!45,
  title={\textsc{Remarque}\ifx\relax#1\relax\relax\else\textmd{ : \itshape #1}\fi}}
% --- boîtes EXERCICE / CORRIGÉ (noms EXACTS, auto-comptées) ---
\newtcolorbox[auto counter]{exo}[1][]{boxbase,colback=primary!4,colframe=primary,
  title={\textsc{Exercice}~\thetcbcounter\ifx\relax#1\relax\relax\else\textmd{ --- \itshape #1}\fi}}
\newtcolorbox[auto counter]{corrige}[1][]{boxbase,colback=excol!4,colframe=excol,
  title={\textsc{Corrigé}~\thetcbcounter\ifx\relax#1\relax\relax\else\textmd{ --- \itshape #1}\fi}}
\newcommand{\vv}[1]{\vec{#1}}\newcommand{\ds}{\displaystyle}
\newcommand{\R}{\mathbb{R}}\newcommand{\N}{\mathbb{N}}\newcommand{\Z}{\mathbb{Z}}
% =========================================================================
\begin{document}

\begin{exo}[où la réflexion totale est-elle possible ?]
On dispose de trois milieux : air ($n=1$), eau ($n=1{,}33$) et verre ($n=1{,}5$).
\begin{enumerate}[label=\alph*)]
  \item Parmi tous les passages possibles entre deux de ces milieux, lesquels peuvent donner lieu à
        une réflexion totale (à condition d'un angle assez grand) ?
  \item Calcule l'angle limite de chacun de ces passages.
\end{enumerate}
\end{exo}

\begin{exo}[le prisme à réflexion totale]
Un rayon entre \emph{perpendiculairement} dans la petite face verticale d'un prisme en verre
($n=1{,}5$) dont la section est un triangle rectangle isocèle ($\ang{45}$–$\ang{45}$–$\ang{90}$),
entouré d'air.
\begin{center}
\begin{tikzpicture}[scale=1.15,>=Stealth]
  \coordinate (A) at (0,0); \coordinate (B) at (0,3); \coordinate (C) at (3,0);
  \fill[defcol!25] (A)--(B)--(C)--cycle;
  \draw[black!70,line width=1pt] (A)--(B)--(C)--cycle;
  \node[font=\scriptsize] at (0.85,0.75){verre $n=1{,}5$};
  \node[above left,font=\scriptsize] at (B){\ang{45}};
  \node[below right,font=\scriptsize] at (C){\ang{45}};
  \draw (0.22,0)--(0.22,0.22)--(0,0.22);
  \coordinate (M) at (0,1.8);
  \draw[->,line width=1pt] (-1.3,1.8)--(M);
  \coordinate (N) at (1.2,1.8);
  \draw[line width=1pt,primary] (M)--(N);
  \fill (N) circle (1pt); \node[above right,font=\scriptsize] at (N){$N$};
  \draw[dash pattern=on 3pt off 2pt,black!60] (N)--($(N)+(0.7,0.7)$);
  \draw[line width=0.5pt] ($(N)+(0.5,0)$) arc (0:45:0.5);
  \node[font=\scriptsize] at (1.82,2.0){\ang{45}};
  \coordinate (P) at (1.2,0);
  \draw[line width=1pt,primary] (N)--(P);
  \draw[->,line width=1pt,primary] (P)--(1.2,-1.2);
  \node[left,font=\scriptsize] at (-1.3,1.95){rayon};
\end{tikzpicture}
\end{center}
\begin{enumerate}[label=\alph*)]
  \item Sous quel angle le rayon frappe-t-il l'hypoténuse (par rapport à sa normale) ?
  \item La réflexion sur l'hypoténuse est-elle totale ? (angle limite verre/air $\approx\ang{41.8}$)
  \item De quel angle le prisme fait-il finalement tourner le rayon ?
\end{enumerate}
\end{exo}

\begin{exo}[un rayon dans une fibre optique]
Un rayon pénètre dans une fibre optique en verre ($n=1{,}5$, entourée d'air) et chemine à
l'intérieur en faisant un angle de $\ang{40}$ \textbf{avec la face d'entrée} (verticale). Il frappe
ensuite la paroi latérale (horizontale) en $P$.
\begin{center}
\begin{tikzpicture}[scale=1.0,>=Stealth]
  \fill[defcol!22] (0,-1.1) rectangle (6.5,1.1);
  \draw[primary,line width=1.4pt] (0,1.1)--(6.5,1.1);
  \draw[primary,line width=1.4pt] (0,-1.1)--(6.5,-1.1);
  \draw[primary,line width=1.4pt] (0,-1.1)--(0,1.1);
  \node[left,font=\scriptsize] at (0,0){face d'entrée};
  \node[font=\scriptsize] at (4.6,-0.5){fibre $n=1{,}5$};
  \coordinate (M) at (0,-0.5);
  \draw[->,line width=1pt] (-1.4,-1.3)--(M);
  \fill (M) circle (1pt); \node[below left,font=\scriptsize] at (M){$M$};
  \draw[dash pattern=on 3pt off 2pt,black!55] (M)--(1.5,-0.5);
  \coordinate (P) at (1.95,1.1);
  \draw[line width=1.1pt,primary] (M)--(P);
  \fill (P) circle (1pt); \node[above,font=\scriptsize] at (P){$P$};
  \draw[dash pattern=on 3pt off 2pt,black!55] (P)--(1.95,0.4);
  \draw[line width=0.5pt] (0.7,-0.5) arc (0:34:0.7);
  \node[font=\scriptsize] at (0.95,-0.28){\ang{40}};
  \draw[line width=1.1pt,primary] (P)--(3.9,-1.1);
\end{tikzpicture}
\end{center}
\begin{enumerate}[label=\alph*)]
  \item Quel est l'angle d'incidence $\hat{\imath}_P$ du rayon sur la paroi en $P$ ?
  \item La réflexion en $P$ est-elle totale ? Le rayon reste-t-il guidé dans la fibre ?
\end{enumerate}
\end{exo}

\begin{exo}[indice minimal pour réfléchir à $\ang{45}$]
Dans un prisme en verre entouré d'air, un rayon frappe une face interne sous un angle d'incidence de
$\ang{45}$. On veut que la réflexion y soit \textbf{totale}.

Détermine l'indice de réfraction \emph{minimal} $n_{\min}$ que doit avoir le verre du prisme.
\end{exo}

\begin{exo}[le prisme immergé]
Un prisme en verre ($n=1{,}5$) est traversé par un rayon qui frappe une face interne sous un angle
d'incidence de $\ang{60}$. On entoure le prisme d'un liquide d'indice $n'$.
\begin{enumerate}[label=\alph*)]
  \item Écris la condition sur $n'$ pour que la réflexion en cette face reste totale.
  \item Quelle est la valeur maximale de $n'$ ? La réflexion reste-t-elle totale dans l'air
        ($n'=1$) ?
\end{enumerate}
\end{exo}

\end{document}
